\section{Executable Evaluation and Cost Analysis}
\label{sec:evidence}

The evaluation asks whether the judges realize the atomic distinctions, which
representation choices determine communication, and which statement layers
dominate a succinct proof. The goal is design insight from an executable model,
not a claim of device feasibility.

\subsection{Artifact and Method}
\label{sec:artifact}

The artifact implements nine complete stacks, one deliberately partial
baseline, five circuits, and four adjacent-condition attacks. A complete stack
carries evidence for every G-condition through its claimed level. The seeded
generator, raw JSON, table generator, numeric claim checker, and adversarial
proof-object tests are included in the anonymous supplement at
\url{https://anonymous.4open.science/r/edge-fusion/}.

We report
$\mathrm{cost}(S,n,m,F)=\langle\mathbf{N}_{P},\mathbf{N}_{J},
B_{\mathrm{ev}},M_P,M_J\rangle$. The first two coordinates count named
evidence-generation and judge primitives. $B_{\mathrm{ev}}$ is the online
transcript beyond the pre-provisioned policy $P$ and claimed output $y$,
including verifier challenges. $M_P$ and $M_J$ are protocol state retained
between records under the declared streaming schedule, not process RSS. If
$B_P$ bytes of policy and setup material are distributed over $R$ runs, the
amortized communication is $B_{\mathrm{ev}}+B_P/R$. The same rule applies to
device keys, policy rosters, anchor keys, and proof verification keys.
The generator populates a canonically serialized, versioned Ed25519/BLS roster;
$H(P)$ binds its exact enrolled keys.

The workload contains $n=60$ scheduled readings from $m=4$ loggers and tests
arithmetic mean, excursion count, and bounded fixed-point mean kinetic
temperature (MKT). It is synthetic scalar sensor-data aggregation, not an
evaluation of tracking, association, heterogeneous-source fusion, or
pharmaceutical prevalence. The stack sweep covers
$n\in\{12,20,60,120,240,480,960\}$. Groth16 experiments use circom 2.2.3,
circomlib 2.0.5, and snarkjs 0.7.5 on an Apple M4 Pro with 24 GB of memory.

\subsection{Separation and the Price of Sampling}
\label{sec:separation}

The adjacent traces alter a record after anchoring (A1), use a credential that
does not match the named device (A2), omit a required identifier (A3-short), or
change the output of a complete authentic run (A4). Replaying each trace against
each stack's own evidence gives 36 complete-stack decisions; 34 are
deterministic and all follow the cumulative definition. These executions are
counterexamples separating adjacent goals, not security proofs for all
adversaries. Separate mutations replace an enrolled key or the accepted roster;
both are rejected.

The sampled V3 stack adds a harder coverage trace, A3-padded: it keeps length
$n$ by replacing the omitted identifier with an authentic duplicate. The
auditor anchors the root, stores its own unpredictable 32-byte challenge, and
checks exactly the $k$ derived policy positions, paths, signatures, and
identifiers. The challenge is not read from the producer's proof object.
A3-short is rejected before sampling because its length violates $P$;
A3-padded is detected only if a defective position is opened.

If $t$ of $n$ positions are defective before the challenge and $k$ distinct
positions are sampled uniformly without replacement, then
\[
 \Pr[\mathrm{miss}]=\delta=\frac{\binom{n-t}{k}}{\binom{n}{k}},\qquad
 \Pr[\mathrm{detect}]=1-\delta,
\]
with $\binom{n-t}{k}=0$ when $k>n-t$. For one defect, detection is $k/n$.
Thus $n=60,k=20$ checks one third of the manifest but still accepts a fixed
single-position defect with probability $2/3$. As an implementation sanity
check, 2,000 fixed challenges detected A1 at 0.338 and A3-padded at 0.317,
consistent with $1/3$. The formula does not cover a producer that learns or
grinds the challenge before committing.

\textbf{Design rule 1:} sampling is attractive only when the policy explicitly
accepts sparse-defect false acceptance; it is not a strong defense against one
strategic omission.

\subsection{Representation, Not the V-Number, Drives Bandwidth}

For the declared serialization, the dominant byte formulas are
\begin{align*}
B_{\mathrm{raw}}(n) &= 144+n(6+64),\\
B_{\mathrm{Ed}}(n) &= 144+n(36+64),\\
B_{\mathrm{BLS}}(n) &= 240+36n,\\
B_{\mathrm{sample}}(n,k) &= 176+k'(100+32\lceil\log_2 n\rceil),
\quad k'=\min(k,n).
\end{align*}
The constants expose the boundary: each anchored stack carries a 32-byte root,
64-byte signature, 16-byte run ID, and 32-byte auditor state; 36 bytes is an
identifier and commitment. Sampled openings also carry a Merkle path. Exact V3 sends the
same full manifest as Ed25519 V2, so it adds judge work but no bytes.

At $n=60$, raw authenticated records and recomputation provide cumulative V4 in
4,344 bytes, while commitment-based Ed25519 V2 uses 6,144 bytes. Here a 6-byte
observation is smaller than its 36-byte manifest entry. With fixed $k=20$, the
sampled stack uses 6,016 bytes at $n=60$; accounting for Merkle-depth steps, it
first becomes smaller than raw transmission at $n=94$ (the first measured point
past that crossover is $n=120$). It buys that reduction with interaction and
$\delta$, not with a stronger primitive.

% generated by artifact/scripts/make_tables.py -- do not edit
\begin{table}[t]
\caption{Decision rules from the executable model. Byte comparisons use
$n=60$ and exclude pre-provisioned $P$ and claimed $y$. The attestation row
is a cost model; all others are executed.}
\label{tab:costvec}
\centering
\scriptsize
\setlength{\tabcolsep}{2.2pt}
\begin{tabular}{@{}L{2.5cm}L{4.2cm}L{4.8cm}@{}}
\toprule
Design choice & Finding & Decision implication \\
\midrule
Raw records vs. commitments & Raw cumulative V4 is 4\,344 B; commitment-based
V2 is 6\,144 B. & Recompute when records are shorter than commitments and
disclosure is acceptable. \\
Exact vs. sampled coverage & Exact V3 reuses the 6\,144-B V2 manifest;
$k=20$ sampling uses 6\,016 B. & Sampling saves only at scale and accepts
an explicit false-acceptance bound. \\
Ed25519 vs. BLS & BLS reduces V2 to 2\,400 B but uses 61 verifier pairings. &
Choose from auditor work and public-verification needs, not tag size alone. \\
Succinct attribution & In-circuit EdDSA expands the mean circuit 17.1$\times$. &
Keep authentication outside the circuit or choose proof-friendly credentials
when the trust model permits. \\
\emph{Modeled:} attestation & A root, external anchor, and report total 208 B
under $\rho_{\mathrm{iso}}$. & Compare the isolated-path and appraisal
assumptions before comparing bytes. \\
\bottomrule
\end{tabular}
\end{table}


The setup comparison also separates encodings. At $n=12$, snarkjs emits 983
bytes of proof and public-signal JSON; the G1 signature, run ID, and auditor
state bring this implementation bundle to 1,095 bytes. Packed raw V4 uses 984
bytes, so this cross-encoding implementation comparison is 111 bytes larger.
Under normalized compressed-binary encoding, the 128-byte proof, 96-byte public fields,
and 112-byte G1 context total 336 bytes---648 bytes smaller than packed raw.
The verification key is 3,290 bytes as JSON or 352 bytes normalized; a first run is
therefore 4,385 bytes in JSON or 688 bytes normalized. The $n=60$ attributed
circuit was compiled but not proven, so we make no comparison there.

Aggregate signatures expose a different trade. BLS reduces V2 evidence from
6,144 to 2,400 bytes but replaces 61 Ed25519 verifications with 61 pairings and
one aggregate-signature verification. The 16-byte homomorphic MAC is smaller
still, but authenticates only a linear aggregate for a designated auditor. It
omits G1 and per-record G2 evidence and is therefore a partial baseline, not V4.

The modeled attestation stack is visually separated in Table~\ref{tab:costvec}.
Its 208-byte transcript contains a 32-byte manifest root, 48-byte run context,
a 64-byte external anchor that supplies G1 before output acceptance, and a
64-byte isolated-path report binding $P$, the root, and $y$. It assumes enrollment, appraisal, and
policy enforcement under $\rho_{\mathrm{iso}}$; no secure element or device was
measured.

\textbf{Design rule 2:} compare complete evidence objects and their trust and
setup boundaries. A higher V-number need not cost more than a lower one, and a
small primitive is not automatically a complete stack.

\subsection{Attribution Dominates the Demonstrated Circuit}

Table~\ref{tab:circuits} decomposes the proof statement. Binding 60
policy-ordered values takes 32,040 constraints. Arithmetic mean adds one;
excursion counting adds 1,261 (3.9\%); and bounded MKT adds 6,591 (20.6\%). The
MKT circuit constrains the declared 0--35\,$^{\circ}$C band at scale $2^{16}$.
Integer divisions and the logarithm round toward $-\infty$ with constrained
remainders; the degree-12 approximation differs from floating point by 1.8~mK
on this trace. These are alternative
G4 operators over the same binding layer.

% generated by artifact/scripts/make_tables.py -- do not edit
\begin{table}[t]
\caption{Constraint waterfall at $n=60$. Operator rows are alternatives over
the binding layer; attribution is added to the arithmetic mean. The attributed
row was compiled but not proven at $n=60$ and still requires external G1.}
\label{tab:circuits}
\centering
\small
\setlength{\tabcolsep}{4pt}
\begin{tabular}{@{}lrrl@{}}
\toprule
Layer & Constraints & Increment & Direct role \\
\midrule
input binding & 32\,040 & --- & manifest binding \\
+ arithmetic mean & 32\,041 & +1 & G4 arithmetic \\
+ excursion count & 33\,301 & +1\,261 & G4 comparison \\
+ bounded MKT & 38\,631 & +6\,591 & G4 approximation \\
+ EdDSA attribution and mean & 549\,392 & +517\,351 & G2--G4 \\
\bottomrule
\end{tabular}
\end{table}


Adding policy-positioned EdDSA attribution to the arithmetic mean raises the
circuit from 32,041 to 549,392 constraints, a factor of 17.1. It constrains
G2--G4 and exposes $c_M$, but an external event must still supply G1. The
$n=60$ attributed circuit was compilation-only. At $n=12$, where the available
setup sufficed, proving consumed a median 2,467 MB of resident memory over three
trials (range 1,812--2,494 MB). This host measurement is evidence against
assuming logger feasibility, not a device benchmark.

\textbf{Design rule 3:} optimizing the fusion formula alone will not make this
proof deployable while conventional in-circuit attribution dominates. Keeping
authentication outside the circuit or using proof-friendly credentials may be
more consequential, when the resulting trust model is acceptable.

\subsection{Threats to Validity}

The artifact is an executable protocol model, not deployable security software.
Fixture keys are publicly derivable, the Groth16 setup is for benchmarking, and
the isolated attestation path is assumed. Primitive and serialized-byte counts
are exact for the code, but retained state is analytical and prover RSS is
host-specific. We report no device latency, energy, radio duty cycle, or
physical-attack result. Commitment openings also require retained source data;
availability and audit-storage cost are outside the cost vector.

One synthetic seed and scalar reductions cannot establish general fusion
workload realism. Order-sensitive tracking or filtering can use the
canonical policy order, but its circuit and protocol costs remain unmeasured.
The sampling expression assumes a defect set fixed before a sound challenge;
32-bit draws use rejection sampling for exact uniformity.
adaptive corruption and open-window acquisition need additional protocols.
Finally, MKT correctness is relative to one bounded numerical specification,
not all temperature bands or implementations.
